\documentclass[12pt,a4paper]{article}
\usepackage[UTF8, fontset=none]{ctex}
\setCJKmainfont{Droid Sans Fallback}
\setCJKsansfont{Droid Sans Fallback}
\setCJKmonofont{Droid Sans Fallback}
\usepackage{amsfonts,amsmath,amsthm,amssymb}
\usepackage{sfmath}
\usepackage{hyperref}
\usepackage{geometry}
\geometry{left=0cm,right=0cm,top=0cm,bottom=0.1cm, footskip=0cm, includefoot=true}
%\usepackage[contents = \thepage, color=yellow, angle=0, scale = 20]{background}
\pagestyle{empty}
%\usepackage{fancyhdr}
%\pagestyle{fancy}
%\fancyhead{}
%\fancyfoot{}
%\cfoot{}
%\lfoot{{\small \color{red} \slshape \thepage 页}}
%\fancyfoot[R]{\thepage}

\usepackage[all]{xy}
\usepackage{color}

\renewenvironment{proof}{{\noindent\it\textbf{Proof}}\quad\it}{\hfill$\square$}

%\theoremstyle{definition}

\newtheorem{thm}[equation]{Theorem}
\newtheorem{cor}[equation]{Corollary}
\newtheorem{lem}[equation]{Lemma}
\newtheorem{prop}[equation]{Proposition}
\newtheorem{defn}[equation]{Definition}
\newtheorem{rem}[equation]{Remark}

\title{last Kervaire}
\author{}

\setcounter{secnumdepth}{6}

\begin{document}

\sffamily
\maketitle
\thispagestyle{empty}

\abstract{}
\tableofcontents

\large

\begin{sloppy}
\clubpenalty=-100
\widowpenalty=-100

\section{谱序列}
\subsection{谱序列}
\paragraph{}
给定abel范畴C，定义C上的具有n分次的谱序列

\begin{rem}
    本项目考虑的filtration都是decreasing
\end{rem}

\subsection{谱序列的收敛性}

\paragraph{}A为一个n-1分次的对象，给定A上的一个filtration，则graded pieces构成一个n分次的对象

\paragraph{}称谱序列E收敛到A，如果存在同构，$E_\infty$到A的graded pieces，之间分次可以相差一个线性变换

\paragraph{}converging 谱序列的范畴：
给定一个$E$到A分次的变换规则，定义相应的谱序列的范畴，态射为一个E之间的同态加上A之间的同态，保持相应的结构

\begin{rem}
    该收敛性为弱意义下的收敛性，本项目不考虑强收敛性等问题
\end{rem}

\paragraph{}A的原始的n-1分次对应的E的分次，我们称之为stem

\paragraph{}A中的filtration对应的E的分次，我们称之为E的filtration次数

\paragraph{}
A中元素$x$具有filtration至少r，考虑$x$在r分次的graded piece中的像，其在E中对应的元素$y$, 我们称y可以detect x

\paragraph{} x被0 detect当且仅当x的filtration至少$r+1$

\paragraph{}
$x,x'$被同一个$y$所detect，当且仅当$x-x'$被更高filtration的东西detect

\subsection{filtration 诱导的谱序列}
\paragraph{}
给定具有filtration的链复形$X$, 我们可以得到相应的谱序列

\paragraph{}
该谱序列收敛到链复形的同调

\paragraph{}
该谱序列具有函子性

\subsection{微分的交叉}
\paragraph{}
给定$E$中的两个微分
$$d:x\xrightarrow{}y$$
$$d':z\xrightarrow{}w$$
我们称$d$被$d'$交叉，如果
\begin{enumerate}
\item $d'$为非平凡的微分
    \item x与z的stem相同，z的filtration严格高于x
    \item y与w的stem相同，w的filtration不高于y
\end{enumerate}

\paragraph{}
2.7给出了更精细的crossing条件

\subsection{extension谱序列}

\paragraph{}
给定一个converging谱序列的同态
$f:(E_1,A_1)\xrightarrow{}(E_2,A_2)$, 则A上的filtration诱导了extension谱序列


\paragraph{} extension微分$d:x\xrightarrow{}y$不被其他微分交叉，当且仅当，对任意被x所detect的$x_0$, $f(x_0)$被y所detect

\paragraph{}
2.10是更精确的版本

\subsection{extension微分和谱序列中微分的交换性}

即2.12及其推论


\section{稳定同伦论}
\subsection{稳定同伦范畴}\label{3}

本节不依赖其他章节

\paragraph{}
稳定同伦范畴$\mathcal{S}$为一个三角范畴

\begin{rem}
    $\mathcal{S}$是指谱形(spectra)的同伦范畴。在形式化中我们直接公理化其三角范畴结构，而不经过模型范畴的构造。
\end{rem}

\paragraph{}
存在球面谱 $\mathbb{S}\in \mathcal{S}$

\paragraph{}
对整数n，
$S^n$为作用n次suspension函子在$\mathbb{S}$上

\paragraph{}
$\pi_nX$为$S^n$到 X 的映射同伦类

\paragraph{}
$\pi_*X$为单分次的abel群

\subsection{上同调}

\paragraph{}
$HF_2$为Eilenberg-MacLane对象

\paragraph{}
$\pi_0(HF_2)=F_2$， 其他的同伦群为0

\paragraph{}
mod 2系数同调定义为
$$H^*(X) = [\Sigma^*X,HF_2]$$

\subsection{Adams谱序列}

\paragraph{}Adams谱序列是$\mathcal{S}$到双分次的converging
谱序列范畴的函子

\paragraph{}
$\pi_*X$上有自然的Adams filtration

\paragraph{}
Adams谱序列收敛到$\pi_*X$, 即Adams filtration的graded piece与Adams 谱序列的$E_\infty$项中相应的群之间存在自然同构

\paragraph{}mod 2的Adams谱序列的$E_2$项中元素都是2阶的 

\paragraph{}Adams filtration可以推广到一般的两个谱形之间的映射

\paragraph{}Adams filtration的刻画：
$f:X\xrightarrow{}Y$具有Adams filtration至少r，如果f能分解为r个映射的复合，且其中每个映射诱导mod 2同调的0映射

\section{synthetic}
\subsection{synthetic谱形}
考虑$HF_2$-synthetic 谱
\paragraph{}Syn 为预三角范畴

\paragraph{} hSyn
为其同伦范畴

\paragraph{}$\Sigma^{m,n}$为一个hSyn上的自等价

\paragraph{}存在自然等价：
$$\Sigma^{m_1,n_1}\Sigma^{m_2,n_2}\cong\Sigma^{m_1+m_2,n_1+n_2}$$


\paragraph{}$\Sigma^{1,0}$为稳定同伦范畴中的suspension函子


\paragraph{}存在自然变换
$$\lambda: \Sigma^{0,-1}\xrightarrow{} Id$$

\paragraph{}定义$X/\lambda^n$为 $\lambda^n|_X$的cofiber

\paragraph{} hSyn可以enriched over $\mathbb{Z}[\lambda]$-模, $Hom(X,Y)$为双分次的abel群$[\Sigma^{*,*}X,Y]$

\paragraph{} $\lambda$的分次是$(0,-1)$

\subsection{Synthetic球}
\paragraph{}存在球对象
$S^{0,0}\in Syn$

\paragraph{}
$S^{m,n} = \Sigma^{m,n}S^{0,0}$

\paragraph{}synthetic同伦群为bigraded abel群
$$\pi_{m,n}(X) = [S^{m,n},X]$$

\paragraph{}
$\pi_{m,n}(X) = \pi_{m+k,n+l}(\Sigma^{k,l}X)$

\paragraph{}$\pi_{*,*}$具有自然的$\mathbb{Z}[\lambda]$-模结构

\subsection{Adams谱序列}

\paragraph{}Adams谱序列是$\mathcal{S}$到3-分次的converging谱序列范畴的函子，分次规则见3.6

\paragraph{}
$\pi_*X$上有自然的Adams filtration

\paragraph{}
Adams谱序列收敛到$\pi_{*,*}X$

\paragraph{}synthetic Adams谱序列是一个$\mathbb{Z}[\lambda]$-模范畴中的谱序列

\subsection{$\nu$函子}
\paragraph{} 存在函子
$$\nu: h\mathcal{S}\xrightarrow{} hSyn$$

\paragraph{}$\nu$保distinguished triangla的条件：3.1

\paragraph{} $\nu$与$\Sigma$不交换，
$$\nu(\Sigma X) \cong \Sigma^{1,1}\nu X$$

\paragraph{}
自然的比较映射
$\Sigma(\nu X) \xrightarrow{} \nu(\Sigma X)$为$\lambda$

\subsection{Synthetic rigidity}
\paragraph{}
$\nu X$的Adams $E_2$项：即3.6

\paragraph{} Adams 与$\lambda$-bockstein:即3.8
\begin{rem}
    3.8(2)中的映射为cofiber
    $$\Sigma^{0,r}\nu X\xrightarrow{\lambda^r} \nu X \xrightarrow{} \nu X/\lambda^r$$
    中的边缘映射
\end{rem}

\paragraph{}$\nu X$的Synthetic Adams $E_\infty$项与$X$的Adams的关系：即3.12

\subsection{Synthetic lift}

\paragraph{}$\nu$作用在Adams filtration k的映射上得到的东西可以除$\lambda^k$, 即Adams filtration等价于$\lambda$-bockstein filtration, 即3.16

\paragraph{} Adams filtration 1 的映射作用$\nu$之后的canonical 除$\lambda$：即3.17

\paragraph{} $\nu$作用在distinguished triangle上的notation： 3.19，3.20


\begin{thebibliography}{99}


\end{thebibliography}

\end{sloppy}
\end{document}
